% Converted from README.md
\paragraph{Detekcija neispravnog argumenta pri alokaciji memorije}\label{detekcija-neispravnog-argumenta-pri-alokaciji-memorije}

Memcheck može da detekuje grešku slanja negativnog argumenta za veličinu alociranog prostora. Program \texttt{06\_fishy\_arguments.c} ilustruje tu grešku. Ukoliko se program pokrene za negativne brojeve neće se ništa desiti jer petlja u tom slučaju neće imati ni jednu iteraciju, pa neće biti ni poziva \texttt{malloc} funkcije. Pozivi sa \texttt{n\ \textless{}\ 27} neće biti problematični - sve će se lepo alocirati i osloboditi. Imaćemo samo prijavljeno upozorenje:

\begin{minted}{text}
(Warning: set address range perms: large range)
\end{minted}

Ovo znači da se velikom bloku memorije menjaju prava pristupa. To upozorenje je namenjeno zapravo najviše Valgrind developerima i možemo ga ignorisati jer će memorija biti alocirana. Ukoliko pokrenemo ponovo Valgrind i za \texttt{n} unesemo \texttt{28}, dobićemo grešku jer smo prekoračili opseg \texttt{int} domena.

\begin{minted}{text}
sada je i = 26 allocated = 1073741824
==5105== Warning: set address range perms: large range [0x59c89028, 0x99c89058) (noaccess)
sada je i = 27 allocated = -2147483648B
==5105== Argument ’size’ of function malloc has a fishy (possibly negative) value: -2147483648
==5105== at 0x483874F: malloc (in /usr/lib/x86_64-linux-gnu/valgrind/vgpreload_memcheck-amd64-linux.so)
==5105== by 0x109219: main (06_fishy_arguments.c:20)
==5105==
Realloc failed!
\end{minted}
